\section{Discussion: Comparative Dynamics Across All Benchmark Models}
\label{sec:discussion}

The empirical findings from our rolling-origin benchmark provide crucial insights into how distinct algorithmic architectures handle highly autocorrelated, noisy, and regime-shifting financial time series over a $100$-day forecasting horizon. While modern data science literature frequently advocates for deep non-linear models, our evaluation across four distinct macroeconomic interest rate regimes ($2009$--$2025$) reveals a more complex reality. 

This section critically evaluates the performance, structural strengths, and architectural constraints of all six benchmarked models: Feedforward Neural Networks (FFN), Regularized Linear Regression (LinReg), Temporal Fusion Transformers (TFT), Gradient Boosted Trees (XGBoost), Naive Random Walk (Naive RW), and Naive Hold.

\subsection{Feedforward Neural Networks (FFN): High Capacity and Feature Sensitivity}
The Feedforward Neural Network achieved the lowest overall Mean Absolute Error (MAE = $7.25$ Bn. \texteuro), representing a $38\%$ reduction in prediction error relative to the Naive RW baseline. The model's hidden architecture ($128 \times 64$ fully connected nodes) successfully captured complex, non-linear relationships between broad macroeconomic indicators and deposit flows—such as non-linear behavioral shifts during negative rate regimes.

However, cross-regime variance analysis reveals that the FFN is highly sensitive to the chosen feature space. When supplied with rich, multi-dimensional macroeconomic inputs (\texttt{nur\_makro}), the network generalized exceptionally well. Conversely, when restricted to single-rate covariates (e.g., \texttt{nur\_zinsen}), performance degraded rapidly, producing large error spikes. This highlights that deep architectures require multi-vector regularizing inputs to prevent hidden layers from memorizing historical baseline trends that fail to persist across future rolling origins.

\subsection{Regularized Linear Regression (LinReg): The Gold Standard for Stability}
While ranking second in absolute MAE ($8.41$ Bn. \texteuro), the Linear Regression model fitted with Ridge ($L_2$) regularization demonstrated the highest out-of-sample coefficient of determination ($R^2 = 0.255$) among all evaluated approaches. 

The structural advantage of $L_2$-regularized linear models in cumulative horizon forecasting lies in their smooth penalty mechanism. When predicting daily volume differences ($\Delta V_t$) over $H = 100$ steps, small daily estimation errors compound linearly or quadratically during final volume reconstruction:
\begin{equation}
    \hat{V}_{t+h} = V_t + \sum_{k=1}^{h} \Delta \hat{V}_{t+k}
\end{equation}
The Ridge penalty shrinks non-essential feature weights toward zero, effectively filtering out high-frequency noise. Consequently, during extreme regime shifts—such as the transition from near-zero rates in 2013 to aggressive tightening in 2025—LinReg maintained low error variance ($\pm 11.61$ Bn. \texteuro), outperforming all complex models in cross-regime stability.

\subsection{Temporal Fusion Transformer (TFT): Architectural Overhead vs. Input Constraints}
The Temporal Fusion Transformer ranked third overall (MAE = $9.33$ Bn. \texteuro, $R^2 = -0.006$). Although TFT is designed specifically for multi-horizon time series forecasting via attention mechanisms and variable selection networks, its full potential was constrained by the strict ex-post evaluation setup.

To prevent look-ahead bias, future values of dynamic covariates were held constant using a naive-hold assumption over the $100$-day horizon. As a result, TFT’s temporal self-attention layers received static input sequences for future steps, neutralizing its primary structural mechanism for dynamic pattern recognition. While TFT successfully outperformed tree-based models and naive baselines, its heavy parameterization added computational overhead without delivering superior accuracy over simpler FFN or Ridge models under static future covariate assumptions.

\subsection{Gradient Boosted Trees (XGBoost): The Extrapolation Trap}
XGBoost exhibited severe structural limitations in this experiment, yielding a negative out-of-sample $R^2$ ($-0.307$) and a higher average MAE ($10.56$ Bn. \texteuro). 

The primary cause of XGBoost's underperformance is the inherent inability of decision trees to extrapolate continuous trends beyond the bounded feature space of the training data. Tree-based algorithms partition the feature space into orthogonal hyper-rectangles, predicting a constant step-value within each leaf node. When forecasting during historical volume highs (such as the $2025$ rate surge), XGBoost cannot project a continuing upward slope; it caps predictions at the maximum values observed in the training set. Furthermore, because XGBoost was trained via direct multi-output regression ($100$ separate tree ensembles), distant forecast steps ($h \to 100$) accumulated significant variance during summation.

\subsection{Naive Baselines: Essential Anchors for Model Validation}
The inclusion of non-parametric naive benchmarks provides crucial context for evaluating model value:
\begin{itemize}
    \item \textbf{Naive Random Walk (Naive RW):} Assumes zero daily change ($\Delta V = 0$), resulting in an average MAE of $11.78$ Bn. \texteuro{} ($R^2 = -0.564$). While competitive during tranquil market conditions (e.g., $2019$), it fails completely during macroeconomic pivots. All primary ML models (FFN, LinReg, TFT, XGBoost) statistically outperformed Naive RW, confirming that ML models extract genuine economic signals.
    \item \textbf{Naive Hold:} Projects the historical average change, resulting in the worst overall performance (MAE = $17.38$ Bn. \texteuro, $R^2 = -2.246$). This confirms that static trend projection is entirely unsuited for multi-step banking deposit horizons.
\end{itemize}

\subsection{Synthesis and Practical Implications}
Table~\ref{tab:complete_model_comparison} summarizes the empirical rankings and architectural trade-offs across all six evaluated models.

\begin{table}[h!]
\centering
\caption{Comprehensive Benchmark Summary across 4 Interest Rate Regimes ($2009$--$2025$)}
\label{tab:complete_model_comparison}
\begin{tabular}{c l l c c c}
\hline
\textbf{Rank} & \textbf{Model} & \textbf{Best Covariate Set} & \textbf{Mean MAE [Bn. \texteuro]} & \textbf{Mean $R^2$} & \textbf{vs. Naive RW} \\ \hline
\#1 & \textbf{FFN} & \texttt{nur\_makro} & \textbf{7.25} $\pm$ 8.84 & 0.181 & $-38\%$ \\
\#2 & \textbf{LinReg (Ridge)} & \texttt{alle} & 8.41 $\pm$ 11.61 & \textbf{0.255} & $-29\%$ \\
\#3 & \textbf{TFT} & \texttt{nur\_makro} & 9.33 $\pm$ 12.87 & $-0.006$ & $-21\%$ \\
\#4 & \textbf{XGBoost} & \texttt{nur\_einlagezins} & 10.56 $\pm$ 11.37 & $-0.307$ & $-10\%$ \\
\#5 & \textbf{Naive RW} & \texttt{naive} & 11.78 $\pm$ 12.81 & $-0.564$ & Baseline \\
\#6 & \textbf{Naive Hold} & \texttt{naive} & 17.38 $\pm$ 10.13 & $-2.246$ & $+48\%$ worse \\ \hline
\end{tabular}
\end{table}

In conclusion, for practical bank asset-liability management (ALM), \textbf{Regularized Linear Regression} offers the most balanced tradeoff between strong explainability, computational efficiency, and high regime-shift stability ($R^2 = 0.255$), while the \textbf{Feedforward Neural Network} provides the highest absolute precision when rich macroeconomic indicators are available.
